\chapter{Giới thiệu}
\label{chap:intro}

Khi quyền của người dùng thay đổi, máy chủ có thể vẫn chấp nhận thông tin xác thực đã cấp trước đó cho thao tác đã mất quyền. Chương này giới thiệu bài toán phát hiện lỗ hổng này, cùng mục tiêu, phạm vi và đóng góp của khóa luận.

\section{Bối cảnh và động lực nghiên cứu}
\label{sec:intro_background}

OWASP API Security Top 10 năm 2023 xếp lỗ hổng phân quyền ở mức đối tượng và mức chức năng vào nhóm rủi ro hàng đầu \cite{owasp2023api}. Khi hệ thống yêu cầu thu hồi quyền có hiệu lực ngay, nguy cơ an ninh phát sinh nếu phiên cũ vẫn giữ được quyền này. Hiện tượng máy chủ tiếp tục chấp nhận thao tác theo quyền đã thu hồi được gọi là lỗ hổng phân quyền có trạng thái (Stateful Broken Access Control).

Lỗi này đã xuất hiện ở các dự án mã nguồn mở. Trình quản lý tải pyLoad mang lỗ hổng \mbox{CVE-2026-41133} với điểm CVSS 8,8: ứng dụng chép vai trò vào phiên lúc đăng nhập và tiếp tục cấp quyền theo bản sao đó sau khi quản trị viên đổi vai trò \cite{cve202641133}. Máy chủ lưu trữ gói phần mềm (artifact registry) artifact-keeper ghi cờ quản trị vào token JWT lúc đăng nhập, nên quản trị viên bị hạ quyền vẫn giữ quyền quản trị tới khi token hết hạn sau 30 phút; nhà phát triển đã vá lỗi này \cite{artifactkeeper1821}. Hai ứng dụng dùng hai cơ chế xác thực khác nhau nhưng có chung một nguyên nhân: quyền được đọc từ bản sao cấp lúc đăng nhập thay vì từ trạng thái hiện tại.

EvoMaster áp dụng kiểm thử dựa trên tìm kiếm để tự động sinh ca kiểm thử cho REST API \cite{arcuri2019evomaster}. Với lỗi phân quyền, mã trạng thái của một yêu cầu HTTP thường không đủ để biết kết quả có đúng chính sách truy cập hay không. Các oracle phân quyền, tức cơ chế phán quyết kết quả kiểm thử, hiện có trong EvoMaster đối chiếu hành vi giữa các tài khoản có quyền khác nhau \cite{sahin2024bola}. Các oracle này xét mỗi lời gọi tại một thời điểm, nên một lời gọi thành công bằng token cũ sau khi tài khoản bị hạ quyền vẫn trông như một lời gọi hợp lệ. Khi nhánh chính của EvoMaster được chạy trên một ứng dụng mang lỗi này, lưu lượng HTTP có bốn lời gọi như vậy nhận mã 200, nhưng không lời gọi nào bị báo lỗi (Mục~\ref{sec:eval_baseline}).

\section{Phát biểu bài toán}
\label{sec:intro_problem}

Trong phạm vi khóa luận, lỗ hổng phân quyền có trạng thái xảy ra khi hệ thống tiếp tục chấp nhận một thao tác phụ thuộc vào quyền đã thu hồi. Thao tác đó được gửi bằng thông tin xác thực cấp trước lần thay đổi quyền. Chẳng hạn, sau khi một tài khoản bị hạ xuống vai trò thấp hơn, máy chủ vẫn chấp nhận token hoặc cookie cấp trước đó cho thao tác không còn được phép. Tùy chính sách thu hồi quyền và cơ chế xác thực của hệ thống, hiện tượng này có thể thuộc điểm yếu CWE-1464 (Improper Session Invalidation), tức hệ thống không vô hiệu hóa phiên khi cần \cite{cwe1464}. Hình~\ref{fig:stateful_authz_flaw} minh họa ba giai đoạn của lỗ hổng: thao tác ghi ban đầu, thao tác hạ quyền và tái sử dụng token cũ.

Khóa luận tập trung tự động phát hiện việc máy chủ tiếp tục chấp nhận quyền đã thu hồi. Hành vi này không thể phán quyết chỉ từ một yêu cầu HTTP riêng lẻ. Phán quyết cần kết quả của cùng thao tác trước và sau khi thay đổi quyền, cùng dấu hiệu cho thấy quyền đã thật sự bị thu hồi.

\begin{figure}[H]
  \centering
  % Hình 1.1 (flowchart, thay sequence diagram fig_1_1_stateful_authz_flaw): vòng đời lỗ hổng phân quyền có trạng thái.
% Vẽ trực tiếp bằng TikZ để là hình vector, cùng font với thân bài.
\begin{tikzpicture}[
  >=Stealth,
  font=\small\linespread{1}\selectfont,
  node distance=8mm and 7mm,
  procbox/.style={rectangle, draw, rounded corners=2pt, fill=gray!10, align=center,
                  text width=4.1cm, minimum height=1.5cm},
  decision/.style={diamond, draw, fill=gray!10, align=center, aspect=2.2, inner sep=1pt},
  safe/.style={rectangle, draw, rounded corners=6pt, align=center, text width=4.1cm, minimum height=1.0cm},
  vuln/.style={rectangle, draw, very thick, rounded corners=6pt, align=center, text width=4.1cm, minimum height=1.0cm},
  lbl/.style={font=\footnotesize}
]
  \node[procbox] (s1) {\textbf{1. Thao tác ghi ban đầu}\\Tài khoản vai trò quản trị gửi thao tác ghi bằng \mbox{token $c$}, nhận 2xx};
  \node[procbox, right=of s1] (s2) {\textbf{2. Hạ quyền}\\Quản trị viên hạ vai trò của tài khoản; token $c$ vẫn còn hạn};
  \node[procbox, right=of s2] (s3) {\textbf{3. Tái sử dụng token cũ}\\Tài khoản gửi lại thao tác ghi bằng token $c$};
  \node[decision, below=of s3] (d) {Máy chủ tra lại\\vai trò hiện hành?};
  \node[safe, below=8mm of d] (ok) {403: từ chối\\(quyền đã thu hồi có hiệu lực)};
  \node[vuln, left=14mm of d] (bad) {2xx: vẫn chấp nhận\\(lỗ hổng phân quyền có trạng thái)};

  \draw[->] (s1) -- (s2);
  \draw[->] (s2) -- (s3);
  \draw[->] (s3) -- (d);
  \draw[->] (d) -- node[lbl, right] {Có} (ok);
  \draw[->] (d) -- node[lbl, above] {Không} (bad);
\end{tikzpicture}

  \caption{Vòng đời lỗ hổng phân quyền có trạng thái}
  \label{fig:stateful_authz_flaw}
\end{figure}

\section{Mục tiêu và phạm vi nghiên cứu}
\label{sec:intro_objective}

Khóa luận mở rộng EvoMaster để tự động phát hiện các vi phạm phân quyền có trạng thái thuộc quan hệ biến hình MR-A1. Phần mở rộng là một oracle mới, gọi là Oracle 970 theo mã lỗi gán cho vi phạm. Phần mở rộng chạy trong pha kiểm thử bảo mật, sau pha tìm kiếm, nên thuật toán tìm kiếm của công cụ được giữ nguyên.

Về phạm vi, khóa luận chỉ xét quan hệ MR-A1, trong đó quyền thực hiện một thao tác được kiểm tra trước và sau khi vai trò của chủ thể bị hạ. Kịch bản kiểm thử MR-A1 được dựng từ đặc tả OpenAPI, cấu hình xác thực, các quan sát HTTP và kho lưu trữ ca kiểm thử của EvoMaster. Phương pháp chỉ áp dụng khi có đủ ba điều kiện: thao tác hạ quyền thực hiện được qua REST API, hệ thống cho phép đăng nhập lại để lấy thông tin xác thực mới sau khi hạ quyền, và hệ thống từ chối thao tác không đủ quyền bằng mã 403. Các tình huống không thỏa các điều kiện này nằm ngoài phạm vi đánh giá. Chẳng hạn, thao tác xóa chủ thể không để lại tài khoản nào để đăng nhập lại.

\section{Đóng góp của khóa luận}
\label{sec:intro_contributions}

Khóa luận có bốn đóng góp chính:

\begin{enumerate}
    \item \textbf{Quan hệ biến hình MR-A1.} Định nghĩa quan hệ MR-A1 để đối chiếu hành vi phân quyền trước và sau thao tác hạ quyền, với thông tin xác thực giữ cố định. Quan hệ được cài đặt bằng bộ ba thao tác gồm thao tác gốc (source), thao tác chuyển trạng thái (transition) và thao tác kiểm chứng (follow-up), cùng điều kiện tiền đề.
    \item \textbf{Tự động nhận diện thao tác chuyển trạng thái.} Xếp hạng thao tác hạ quyền từ đặc tả OpenAPI bằng mô hình Naive Bayes; cùng mô hình này chọn trường vai trò và trường định danh tài khoản nạn nhân.
    \item \textbf{Kiểm chứng bằng bước đăng nhập lại.} Tài khoản nạn nhân lấy token mới và chạy lại chuỗi lời gọi của mình; bộ ba bị loại nếu token mới không bị từ chối.
    \item \textbf{Đánh giá thực nghiệm.} Đánh giá Oracle 970 trên bốn ứng dụng. Hai ứng dụng reports-api và vampi-seeded mang lỗi tự gieo. Hai ứng dụng mã nguồn mở mang lỗi thật: pyLoad với lỗ hổng \mbox{CVE-2026-41133}, và artifact-keeper với một lỗi đã được nhà phát triển vá. Trên các biến thể có lỗ hổng, Oracle 970 phát hiện vi phạm ở 34/50 lượt, gồm 5/5 lượt trên pyLoad và 4/5 lượt trên artifact-keeper. Trong 48 lượt chạy trên biến thể đã vá, 30 lượt dựng được bộ ba và Oracle 970 không báo động giả ở lượt nào.
\end{enumerate}

\section{Bố cục khóa luận}
\label{sec:intro_organization}

Sau Chương~\ref{chap:background} về cơ sở lý thuyết và các công trình liên quan, Chương~\ref{chap:method} định nghĩa quan hệ MR-A1, cách dựng kịch bản kiểm thử và cách nhận diện thao tác hạ quyền. Cách tích hợp Oracle 970 vào EvoMaster, các ứng dụng mục tiêu và thiết lập thực nghiệm nằm ở Chương~\ref{chap:impl}. Chương~\ref{chap:eval} đánh giá khả năng phát hiện vi phạm, khả năng xếp hạng thao tác chuyển trạng thái và các giới hạn quan sát được, còn Chương~\ref{chap:conclusion} tổng kết kết quả và đề xuất hướng phát triển.
